\documentclass[aspectratio=169]{beamer}

\usepackage[ngerman]{babel}
\usepackage[T1]{fontenc}
\usepackage[utf8]{inputenc}

\usetheme{CambridgeUS}
\usecolortheme{default}

\usepackage{graphicx}
\usepackage{amsmath}
\usepackage{siunitx}
\usepackage{booktabs}

%\usepackage{twemojis}
\title[Kurztitel]{Kick-off Präsentation}
\subtitle{RO45 Robotik Projekt 3 – Automatische Sortieranlage}
\author{NimSort42}
\institute{Hochschule Kempten}
\date{\today}

\setbeamertemplate{section in toc}[ball unnumbered]
\setbeamertemplate{subsection in toc}[ball unnumbered]
\setbeamercolor{structure}{fg=black}


\begin{document}
	
	\begin{frame}
		\titlepage
	\end{frame}
	
	\begin{frame}{Inhalt}
		\tableofcontents
	\end{frame}
	
	% --------------------
	\section{Team \& Projekt}
	
	\begin{frame}{Projektteam \& Projektvorstellung}
		\textbf{Projektteam:}
		\begin{itemize}
			\item Bachhuber Yannick
			\item Keppler Benjamin
			\item Moser Louis
		\end{itemize}
		
		\vspace{6mm}
		
		\textbf{Projektvorstellung:}
		\begin{itemize}
			\item Entwicklung einer Software für eine automatische Sortieranlage
			\item Unterscheidung von Katzen und magischen Einhörnern
			\item Zielgerichtetes Sortieren über einen Portalroboter
		\end{itemize}
	\end{frame}
	
	\section{Projektmanagement}
	
	\begin{frame}{Meilensteine}
		\begin{itemize}
			\item Notwendige Koordinatensysteme festgelegt (30.03.2026)
			\item Code-Grundlagen realisieren (13.04.2026)
			\item Prädizierte Positionen im Weltkoordinatensystem ausgeben (27.04.2026)
			\item Regelung auf einen Punkt im Weltkoordinatensystem (11.05.2026)
			\item Prozesslogik ist in Python programmiert und getestet (01.06.2026)
			\item Abschlusspräsentationen (29.06.2026)
		\end{itemize}
	\end{frame}
	
	\begin{frame}{Risikomanagement}
		\textbf{Mögliche Risiken:}
		\begin{itemize}
			\item \begin{tabular}[t]{@{}p{7.5cm} l@{}}
				Hardware steht nicht rechtzeitig zur Verfügung & $\Rightarrow$ Zeitverzug
			\end{tabular}
			\item \begin{tabular}[t]{@{}p{7.5cm} l@{}}
				Ausfall eines Teammitglieds & $\Rightarrow$ Mehrbelastung des Teams
			\end{tabular}
			\item \begin{tabular}[t]{@{}p{7.5cm} l@{}}
				Professor fällt aus & $\Rightarrow$ Projekt kann nicht bewertet werden
			\end{tabular}
			\item \begin{tabular}[t]{@{}p{7.5cm} l@{}}
				Eingeschränkte Hardwareverfügbarkeit & $\Rightarrow$ Weniger Tests
			\end{tabular}
		\end{itemize}
		\textbf{Umgang:}\\
		\begin{itemize}
			\item Akzeptanz
			\item Puffer eingeplant (4 Wochen)
		\end{itemize} 
	\end{frame}
	
	\begin{frame}{Stakeholder}
		\begin{itemize}
			\item \begin{tabular}[t]{@{}p{3cm} l@{}}
				Teammitglieder & $\Rightarrow$ Bestehen des Moduls
			\end{tabular}
			\item \begin{tabular}[t]{@{}p{3cm} l@{}}
				Professor & $\Rightarrow$ Durchführung und Bewertung des Projekts
			\end{tabular}
			\item \begin{tabular}[t]{@{}p{3cm} l@{}}
				Kunde (Haribo) & $\Rightarrow$ Zuverlässige Sortieranlage
			\end{tabular}
		\end{itemize}
	\end{frame}
	
	
	\begin{frame}{Hauptverantwortlichkeiten}
		\begin{tabular}{@{}l l@{}}
			\textbf{Bachhuber Yannick} & $\Rightarrow$ Prozesslogik, Initialisierung \\ 
			\textbf{Keppler Benjamin} & $\Rightarrow$ Regelung, Achsensteuerung \\
			\textbf{Moser Louis} & $\Rightarrow$ Vision-Pipeline, Objekterkennung \\
		\end{tabular}
	\end{frame}
	
	% --------------------
	\section{Technische Realisierung}
	
	\begin{frame}{Hardware}
		\begin{itemize}
			\item Portalroboter mit Sauggreifer (Arbeitsraum: $\SI{0.1}{m}^3$)
			\item Förderband mit variabler Geschwindigkeit ($0$–$\SI{0.05}{m/s}$)
			\item Arduino Mega mit Stepper Shield zur Achsansteuerung
			\item ELP USB-Kamera (2.0 MP) zur Objekterkennung
			\item Zwei Sortierkisten (Katze / Einhorn) und eine Restkiste
		\end{itemize}
	\end{frame}
	
	\begin{frame}{Software}
		\begin{itemize}
			\item Betriebssystem: Ubuntu 22.04
			\item Framework: ROS2 Humble Hawksbill
			\item OpenCV/scikit-learn in Python
			
		\end{itemize}
	\end{frame}
	
	\begin{frame}{Architektur-Achsen}
		\begin{center}
			\includegraphics[width=0.8\textwidth]{arch1.png}
		\end{center}
	\end{frame}
	
	\begin{frame}{Architektur-Prozesslogik}
		\begin{center}
			\includegraphics[width=0.85\textwidth]{arch2.png}
		\end{center}
	\end{frame}	
	
	\begin{frame}{Architektur-Kamera}
		\begin{center}
			\includegraphics[width=0.7\textwidth]{arch3.png}
		\end{center}
	\end{frame}	
	
	\begin{frame}
		\begin{columns}
			\begin{column}{0.6\textwidth}
				\includegraphics[width=0.7\textwidth]{logo_nimsort42.png}
			\end{column}
			\begin{column}{0.4\textwidth}
				\textbf{So leicht haben Sie noch nie sortiert! Wie aus Zauberhand.}\\
				\vspace{5mm}
				\textbf{Projektteam:}\\
				Bachhuber Yannick\\
				Keppler Benjamin\\
				Moser Louis\\
				\vspace{5mm}
				\textbf{Voraussichtliches Projektende:}\\ 30.03.2026
				
			\end{column}
		\end{columns}
	\end{frame}
	
\end{document}